% Extracto íntegro por residencia del propietario científico definitivo de masas.
\section{Dominio interno y problema de realización}

La Ley General de Masas no parte de una tabla de nombres ni de una lista de
valores experimentales. Su dominio interno es la cadena finita
\[
 \mathcal R_{324}\longrightarrow \mathcal C_{81}
 \longrightarrow \mathcal F_{56}\longrightarrow \mathcal M_{13},
\]
donde \(\mathcal R_{324}\) contiene las \(4\) orientaciones de cada una de las
\(81\) celdas, \(\mathcal F_{56}\) reúne historias de ruta equivalentes y
\(\mathcal M_{13}\) conserva las multisecciones operatorias. El sector nulo
se bifurca antes de aplicar el carácter positivo. Por tanto, ni una tabla
histórica abreviada ni una taxonomía de nombres físicos puede sustituir este
dominio.

Sea \(\mathcal P\) el conjunto de estados físicos individualizados. Un estado
elemental se representa mediante el descriptor
\[
 X=(\mathfrak c,\rho,g,f,q,\epsilon,\varsigma,\mathfrak a),
\]
que conserva clase de representación \(\mathfrak c\), representación interna
\(\rho\), generación \(g\), sabor \(f\), carga \(q\), orientación de hoja
\(\epsilon\), sector estadístico \(\varsigma\) y datos de acoplamiento
\(\mathfrak a\). Un estado compuesto añade constituyentes ordenados, árbol de
composición, frontera común y números cuánticos colectivos. Una resonancia
añade los canales abiertos y la prescripción causal del resolvente.

La realización física debe responder a tres preguntas distintas:

\begin{enumerate}
\item qué fibra, órbita o multisección del atlas corresponde al descriptor;
\item qué operador positivo publica la masa sobre esa fibra;
\item qué acoplamiento produce un escalar, un multiplete, una distribución o
      un polo complejo.
\end{enumerate}

Confundir estas preguntas introduce un selector retrospectivo. La respuesta
correcta se obtiene mediante una relación natural, un operador de fibra y una
compresión espectral, en ese orden.

\section{Producto fibrado entre descriptores y rutas}

Sea \(\mathcal D\) el espacio finito de datos internos ya construidos antes del
contraste: representación, orientación, memoria, frontera, carga, sabor,
color, generación y sector estadístico. Denotemos por
\[
 d_{\mathcal P}:\mathcal P\longrightarrow\mathcal D,
 \qquad
 d_{\mathcal R}:\mathcal R_{324}\longrightarrow\mathcal D
\]
los mapas de descriptor físico y descriptor genealógico. Se define el objeto
de realización por el producto fibrado
\[
 \mathcal J
 =\mathcal P\times_{\mathcal D}\mathcal R_{324}
 =\{(X,\gamma):d_{\mathcal P}(X)=d_{\mathcal R}(\gamma)\}.
\]
La fibra de una especie es
\[
\mathscr S(X)=\{\gamma\in\mathcal R_{324}:(X,\gamma)\in\mathcal J\}.
\]
Un descriptor (X) se denomina realizable por el atlas cuando
\[
 d_{\mathcal P}(X)\in\operatorname{im}d_{\mathcal R}.
\]
Esta condición es equivalente a \(\mathscr S(X)\ne\varnothing\). No se deduce
de la mera escritura del producto fibrado: constituye la condición exacta de
existencia que debe verificarse para cada descriptor físico.
El valor de \(\mathscr S\) puede ser una ruta, una órbita o una multisección.
No se fuerza una sección puntual cuando la simetría interna exige conservar
una multiplicidad.

\begin{definition}[Relación de realización admisible]
La relación \(\mathscr S\) es admisible cuando satisface simultáneamente:
\begin{enumerate}
\item depende únicamente de datos construidos antes de abrir la tabla externa;
\item conmuta con la conjugación partícula--antipartícula;
\item es equivariante bajo la inversión celular y las acciones internas de
      calibre;
\item es compatible con refinamiento, retorno nonádico y lectura
      dodecafásica;
\item conserva toda multiplicidad no eliminada por un acoplamiento físico;
\item permanece invariante al permutar o reemplazar los valores externos de
      masa y anchura.
\end{enumerate}
\end{definition}

\begin{theorem}[Realización prospectiva por producto fibrado]
Si \(d_{\mathcal P}\) y \(d_{\mathcal R}\) son equivariantes respecto de una
simetría \(G\) y no dependen de magnitudes externas, entonces \(\mathcal J\)
es un \(G\)-subobjeto de \(\mathcal P\times\mathcal R_{324}\). Para cada
descriptor realizable \(X\), la fibra no vacía \(\mathscr S(X)\) es
\(G_X\)-estable y la asignación \(X\mapsto\mathscr S(X)\) es independiente del
resultado que posteriormente se contraste.
\end{theorem}

\begin{proof}
Para \((X,\gamma)\in\mathcal J\) se tiene
\(d_{\mathcal P}(X)=d_{\mathcal R}(\gamma)\). La equivariancia da
\[
 d_{\mathcal P}(gX)=g\,d_{\mathcal P}(X)
 =g\,d_{\mathcal R}(\gamma)=d_{\mathcal R}(g\gamma),
\]
y, por consiguiente, \((gX,g\gamma)\in\mathcal J\). Si \(g\in G_X\), la
segunda componente permanece en \(\mathscr S(X)\). La independencia de los
valores externos es inmediata porque ninguno de ellos pertenece al dominio
de los dos mapas del producto fibrado.
\end{proof}

El electrón constituye la excepción central. La inversión celular posee un
único punto fijo, la celda \((5,5)\), y por ello su fibra física es de rango
uno. Las restantes multisecciones estables no admiten, en general, un selector
puntual equivariante. Esta diferencia no es una carencia del atlas: es la
distinción matemática entre una sección central y una especie realizada como
órbita o multisección.

\section{Clasificación espectral intrínseca de las especies}

Sea \(\mathcal A_{\rm rut}\) el álgebra conmutativa finita generada por los
observables internos de ruta: firma \(\mathbb Z^5\), carácter, cronología,
\(K_D\), \(K_\Omega\), orientación, memoria y lectores compatibles. En una
fibra \(F\), se define
\[
 \gamma\sim_{\rm sp}\eta
 \quad\Longleftrightarrow\quad
 a(\gamma)=a(\eta)\quad\text{para todo }a\in\mathcal A_{\rm rut}.
\]
Para cada clase \(\lambda\in F/\!\sim_{\rm sp}\),
\[
 P_\lambda=\sum_{\gamma\in\lambda}
 |\gamma\rangle\langle\gamma|
\]
es un proyector espectral; los \(P_\lambda\) son ortogonales y
\(\sum_\lambda P_\lambda=I_F\).

\begin{theorem}[Máxima individualización interna]
Todo clasificador construido exclusivamente con observables de
\(\mathcal A_{\rm rut}\) factoriza de manera única por el cociente
\(F/\!\sim_{\rm sp}\). En consecuencia, dos descriptores con firma completa
distinta pertenecen a sectores distinguibles; dos descriptores que coinciden
para todos los generadores del álgebra no pueden separarse mediante una regla
construida con esos mismos observables.
\end{theorem}

\begin{proof}
Los caracteres de la álgebra conmutativa finita separan exactamente su
espectro conjunto. Un clasificador funcional de sus generadores es constante
en cada fibra del mapa espectral conjunto y, por ello, factoriza por el
cociente. Si dos rutas tienen el mismo carácter sobre toda la álgebra, ningún
elemento de ella puede separarlas.
\end{proof}

Este teorema cierra el límite lógico de la clasificación. Cuando una
representación física exige distinguir dos estados pertenecientes a una misma
clase espectral, debe exhibir un observable adicional o un acoplamiento que
amplíe \(\mathcal A_{\rm rut}\). La ausencia de tal observable no autoriza a
usar la masa externa como etiqueta selectora.

En el sector leptónico cargado, los descriptores distinguen la fibra central,
la multisección de profundidad \(G_2\) y la multisección de profundidad
\(G_4\). La primera produce la línea electrónica única; las otras dos
producen los operadores asociados a las realizaciones muónica y tauónica. Si
el acoplamiento preserva más de un bloque, la predicción correcta es su
espectro; la etiqueta física no autoriza a escoger uno de sus autovalores por
proximidad numérica.


